\item Un cohete modelo es lanzado verticalmente hacia arriba partiendo del reposo. Su aceleración durante los tres primeros segundos es $a(t) = 60t$, tiempo en el cual el combustible se agota y el cohete se convierte en un cuerpo en caída libre. Al cabo de 17 s, se abre el paracaídas del cohete y la velocidad (hacia abajo) disminuye linealmente a razón de $-18 \text{ pies/s}$ en 5 s. Luego el cohete se posa en el suelo a esa velocidad.
\begin{enumerate}
    \item[(a)] Determine la función posición $s$ y la función velocidad $v$ (para todos los valores del tiempo $t$). Trace las gráficas de $s$ y $v$.
    \item[(b)] ¿En qué momento alcanza el cohete su altura máxima y cuál es esa altura?
    \item[(c)] ¿En qué momento aterriza el cohete?
\end{enumerate}
